\begin{question}
\begin{example}[半圆链条的张力与虚功]
均匀柔软链条质量 $m$、长 $\pi R$，沿光滑水平圆柱的上半圆静置，两端均在圆心等高处且自由。求圆顶处的张力，并用虚功计算作核验。
\end{example}
\fig{c2e12q}{两个自由端之间的上半圆链条}{\begin{tikzpicture}[scale=1.3]
\draw[world](0,0)circle(1.5);\draw[very thick,blue](1.5,0)arc(0:180:1.5);
\node[right]at(1.5,0){$A$};\node[left]at(-1.5,0){$C$};\node[above]at(0,1.5){$B$};
\draw[guide](-1.5,0)--(1.5,0);\draw[guide](0,0)--(40:1.5);
\draw(.4,0)arc(0:40:.4)node[midway,right]{$\theta$};
\end{tikzpicture}}
\end{question}
\begin{solution}\textbf{确定变量与基本关系}\par
图 \ref{c2e12q} 的整条链长与四分之一圆链条题不同，线密度为 $\lambda=m/(\pi R)$。沿弧仍有 $\dd T=\lambda gR\cos\theta\dd\theta$，由右端 $T(0)=0$ 得
$T(\theta)=\lambda gR\sin\theta$，因此
\[
 \boxed{T_B=\lambda gR=\frac{mg}{\pi}}.
\]
上式来自 $T(\theta)-T(0)=\lambda gR\int_0^\theta\cos\phi\dd\phi=\lambda gR\sin\theta$。在圆顶 $\theta=\pi/2$ 时，$\sin\theta=1$，代入 $\lambda=m/(\pi R)$ 得 $T_B=mg/\pi$；在左端 $\theta=\pi$ 时，张力重新为零，与另一自由端条件一致。
\fig{c2e12s}{截取右半段，作沿圆周的虚位移}{\begin{tikzpicture}[scale=1.3]
\draw[world](1.5,0)arc(0:90:1.5);\draw[guide](0,0)--(1.5,0);\draw[guide](0,0)--(50:1.5);
\draw[force](0,1.5)--(-.9,1.5)node[left]{$T_B$};
\draw[flow](50:1.5)--++(140:.65)node[above]{$\delta s$};
\draw[guide](50:1.5)--++(0,.418)node[right]{$\delta h=\delta s\cos\theta$};
\end{tikzpicture}}

\par\medskip\Needspace{3\baselineskip}\noindent\textbf{图示分析与后续计算}\par
在图 \ref{c2e12s} 中截取 $AB$ 段，圆顶的另一半链对它施加张力 $T_B$。令该段作沿圆周逆时针的虚位移 $\delta s$，各微段升高 $\delta h=\delta s\cos\theta$，支持力始终与位移垂直。

平衡时拉力功等于重力势能增量：
\[
 T_B\delta s
 =\int_0^{\pi/2}\lambda R\dd\theta\,
 g\,\delta s\cos\theta
 =\lambda Rg\,\delta s.
\]
积分中的 $\lambda R\dd\theta$ 是微段质量，$g\delta s\cos\theta$ 是单位质量的势能增量。$\int_0^{\pi/2}\cos\theta\dd\theta=1$，再约去共同的非零虚位移 $\delta s$，得到 $T_B=\lambda Rg=mg/\pi$。不能直接用整条链重量的一半作为顶端张力，因为沿圆周各位置的切向重力分量不同。本题只讨论所给对称静置状态，微小不对称滑动的稳定性需另行分析。

\end{solution}
